\documentclass[10pt]{article}
\usepackage[T1]{fontenc}
\usepackage[utf8]{inputenc}
\usepackage[french]{babel}
\usepackage{amsmath,amssymb,mathrsfs,mathpazo}
\usepackage{geometry}
\geometry{a4paper,margin=2cm}
\usepackage{graphicx,array,tabularx,tabularray,colortbl,multirow,multicol,booktabs,enumitem,pifont,ulem,fancybox,diagbox,cancel}
\usepackage{pstricks,pstricks-add,pst-plot,pst-tree,pst-node,pst-text,pst-func,pst-math,pst-bspline,pst-3dplot,pst-eucl,pst-solides3d}
\let\pscalculate\undefined
\usepackage{pst-bezier}
\usepackage{tikz,pgfplots,tkz-tab}
\usetikzlibrary{arrows,arrows.meta,patterns,calc,positioning,trees,decorations.markings,shapes.geometric,shapes.misc,shapes.arrows,intersections,decorations.pathmorphing,decorations.pathreplacing,backgrounds,shadows,babel}
\pgfplotsset{compat=1.18}
\usepgfplotslibrary{fillbetween}
\usepackage[np]{numprint}
\usepackage{listings,makecell,pas-tableur,twoopt}
\let\barre\undefined
\usepackage[pstricks]{tablvar}
\definecolor{darkgreen}{RGB}{0,100,0}
\providecommand{\red}{\color{red}}
\providecommand{\blue}{\color{blue}}
\providecommand{\green}{\color{green}}
\providecommand{\black}{\color{black}}
\providecommand{\white}{\color{white}}
\providecommand{\gray}{\color{gray}}
\providecommand{\cyan}{\color{cyan}}
\providecommand{\magenta}{\color{magenta}}
\providecommand{\yellow}{\color{yellow}}
\providecommand{\violet}{\color{violet}}
\providecommand{\orange}{\color{orange}}
\providecommand{\euro}{\texteuro}
\providecommand{\EUR}{\texteuro}
\providecommand{\ieme}{\textsuperscript{e}}
\providecommand{\degres}{\ensuremath{^\circ}}
\providecommand{\degree}{\ensuremath{^\circ}}
\providecommand{\vv}[1]{\overrightarrow{#1}}
\providecommand{\Coord}[1]{#1}
\providecommand{\decosix}{}
\providecommand{\starredbullet}{\ensuremath{\bullet}}
\pagestyle{empty}
\setlength{\parindent}{0pt}
\usepackage{adjustbox}\DeclareUnicodeCharacter{2236}{:}\DeclareUnicodeCharacter{2190}{\ensuremath{\leftarrow}}
\providecommand{\R}{\mathbb{R}}
\providecommand{\N}{\mathbb{N}}
\providecommand{\Z}{\mathbb{Z}}
\providecommand{\Q}{\mathbb{Q}}
\providecommand{\D}{\mathbb{D}}
\providecommand{\C}{\mathbb{C}}
\providecommand{\e}{{\mathrm{e}}}
\providecommand{\vect}[1]{\overrightarrow{#1}}
\providecommand{\vectt}[1]{\overrightarrow{\mathrm{#1}}}
\providecommand{\vvt}[1]{\overrightarrow{\mathrm{#1}}}
\providecommand{\np}[2][]{{#2~#1}}
\providecommand{\vv}[1]{{\overrightarrow{#1}}}
\renewcommand{\d}{{\mathrm{d}}}
\providecommand{\degres}{{^\circ}}
\providecommand{\degre}{{^\circ}}
\providecommand{\blue}{\color{blue}}
\providecommand{\red}{\color{red}}
\providecommand{\enskip}{\quad}
\providecommand{\llbracket}{\lbrack\!\lbrack}
\providecommand{\rrbracket}{\rbrack\!\rbrack}
\providecommand{\euro}{{\text{€}}}
\providecommand{\no}{{\text{n°}}}
\providecommand{\interval}[2]{{[#1;#2]}}
\providecommand{\intervFO}[2]{{[#1;#2[}}
\providecommand{\dotfill}{\ldots}

\providecommand{\ds}{\displaystyle}
\providecommand{\textding}[1]{#1}
\providecommand{\Oij}{\(\mathrm{O};\vec{\imath},\vec{\jmath}\)}
\providecommand{\Oijk}{\(\mathrm{O};\vec{\imath},\vec{\jmath},\vec{k}\)}
\providecommand{\Ouv}{\(\mathrm{O};\vec{u},\vec{v}\)}
\providecommand{\intervFF}[2]{\ensuremath{[#1;#2]}}
\renewcommand{\euro}{\ifmmode\text{\texteuro}\else\texteuro\fi}
\DeclareUnicodeCharacter{00B0}{\ifmmode^{\circ}\else\textdegree\fi}
\DeclareUnicodeCharacter{20AC}{\ifmmode\text{\texteuro}\else\texteuro\fi}
% Énoncés : Éducation nationale. Transcription et corrigés : APMEP, Denis Vergès et contributeurs. Découpage et classement : CoopMaths, Sébastien Lozano, Éric Elter et contributeurs. Adaptation de la présentation : Maths Explorer.
% CC BY-SA : https://coopmaths.fr/www/about/
% https://coopmaths.fr/alea/static/bac/2025/tex/bac_2025_06_sujet2_asie_3_cor.tex
\begin{document}
\noindent\textit{APMEP / CoopMaths. CC BY-SA. Présentation adaptée par Maths Explorer.}\par\medskip
\medskip

\textbf{Affirmation 1 :} 

D'après la représentation paramétrique donnée le vecteur $\overrightarrow{u} \begin{pmatrix}  \frac{1}{3} \\  \frac{1}{2}  \\1 \end{pmatrix}$ est un vecteur directeur de $(d)$ et d'après l'équation cartésienne de $P$ , $\overrightarrow{n} \begin{pmatrix}  2 \\  3  \\6 \end{pmatrix}$  est un vecteur normal à $P$.

$\overrightarrow{n}=6\,\overrightarrow{u}$ donc les vecteurs $\overrightarrow{n}$ et $\overrightarrow{u}$ sont colinéaires donc $(d)$ est orthogonale à $P$.

$2x_H + 3y_H + 6z_H - 6 = 2 \times (-6)+3 \times 2 + 6 \times 2 -6=-6 \neq 0 $ donc $H \notin P$ donc \textbf{ l'affirmation 1 est fausse}.
 

\textbf{Affirmation 2 :}  

$\overrightarrow{CH} \begin{pmatrix} -9 \\ 2 \\ 2 \end{pmatrix}$ et $\overrightarrow{CD} \begin{pmatrix} -3 \\2\\0 \end{pmatrix}$ donc $\overrightarrow{CH} \cdot \overrightarrow{CD}=-9 \times (-3) + 2\times 2 + 2 \times 0=31$

or par définition $\overrightarrow{CH} \cdot \overrightarrow{CD} = \| \overrightarrow{CH} \| \times \| \overrightarrow{CD} \| \times \cos \left( \, \widehat{DCH}\,\right)$

on a  $\| \overrightarrow{CH} \| = \sqrt{(-9)^2+2^2+2^2}=\sqrt{89}$ et $\| \overrightarrow{CD} \| = \sqrt{(-3)^2+2^2+0^2}=\sqrt{13}$ 

donc $\cos \left( \,\widehat{DCH}\,\right) = \dfrac{\overrightarrow{CH} \cdot \overrightarrow{CD} }{ \| \overrightarrow{CH} \| \times \| \overrightarrow{CD} \|}= \dfrac{31}{\sqrt{89} \times \sqrt{13}} $ 

donc l'angle $\widehat{DCH}$ mesure $\arccos \left( \dfrac{31}{\sqrt{89} \times \sqrt{13}} \right)\approx 24^\circ$ donc \textbf{l'affirmation 2 est vraie}.
 
 

\textbf{Affirmation 3 :} $\overrightarrow{n} \begin{pmatrix}  2 \\  3  \\6 \end{pmatrix}$  est un vecteur normal à $P$ ; $\overrightarrow{n'} \begin{pmatrix}  1 \\  -2  \\3 \end{pmatrix}$  est un vecteur normal à $P'$.

Les vecteurs $\overrightarrow{n}$ et $\overrightarrow{n'}$ ne sont pas colinéaires donc les plan $P$ et $P'$ ne sont pas parallèles et sont donc sécants, leur intersection est une droite.

A$(\,3\,;\,0\,;\,0) \in \Delta$ ( obtenu avec $t=0$) 

$2x_A + 3y_A + 6z_A - 6 = 2 \times 3+3 \times0 + 6 \times 0 -6=  0 $ donc $A \in P$  et $x_A - 2y_A + 3z_A - 3 =3 - 2\times0 + 3\times0 - 3=0$ donc $A \in P'$

Donc l'intersection des plans \( P \) et \( P' \)  est une droite contenant A. Pour déterminer si cette droite est $\Delta$, soit on prend un autre point, soit on montre qu'un vecteur directeur de $\Delta$ est orthogonal à $\overrightarrow{n}$ et $\overrightarrow{n'}$.

B$(\,0\,;\,0\,;\,1) \in \Delta$ ( obtenu avec $t=1$) 

$2x_B + 3y_B + 6z_B - 6 = 2 \times 0+3 \times0 + 6 \times 1 -6=  0 $ donc $B \in P$  et $x_B - 2y_B + 3z_B - 3 =0 - 2\times0 + 3\times1 - 3=0$ donc $B\in P'$ 
donc \textbf{l'affirmation 3 est vraie}.

 

\textbf{Affirmation 4 :} $\overrightarrow{CD} \begin{pmatrix} -3 \\2\\0 \end{pmatrix}$ donc une représentation paramétrique de la droite (CD) est représentation paramétrique : $ \left\{
    \begin{array}{l}
    x = 3-3s\\
    y = 2s \\
    z = 0
    \end{array},\quad s \in \mathbb{R}
    \right.$

pour $s=\dfrac{31}{13}$ on retrouve les coordonnées de $J$ donc $J \in$ (CD).

de plus $\overrightarrow{HJ} \begin{pmatrix} \frac{24}{13} \\ \frac{36}{13} \\ -2   \end{pmatrix}$ donc  $\overrightarrow{HJ} \cdot \overrightarrow{CD} = \dfrac{24}{13} \times (-3) + \dfrac{36}{13} \times 2 + (-2) \times 0 = 0$

donc les vecteurs $\overrightarrow{HJ}$ et $\overrightarrow{CD}$ sont orthogonaux donc  \( J \) est bien le projeté orthogonal de \( H \) sur \( (CD) \) donc \textbf{l'affirmation 4 est vraie}.

\bigskip

\end{document}
